\documentclass[11pt]{article}
\usepackage{amsmath,amssymb,amsthm}
\newtheorem{theorem}{Theorem}
\begin{document}
\title{Exact Two-Fraction Identities on the Full Boolean Cube}
\author{}
\date{November 24, 2026}
\maketitle

\paragraph{Status and scope.}
The classification below is a complete written proof with exact rational
checks, not a Lean theorem. It covers all pairs of original admissible
denominators in dimension at least four, including mixed orientations
and constant denominators. It classifies exact score identities, not
strict sign equivalence. It gives no bound for general sums of many
heads and does not establish either accepted asymptotic target.

\begin{theorem}
Let $n\geq4$. Let $A,B,C,D$ be real affine forms on $\mathbb R^n$.
Suppose $B,D$ are positive at all Boolean vertices, and each denominator
is either constant or has a nonzero coefficient in every variable.
Suppose $A/B=C/D$ at every Boolean vertex. Then the polynomial
$q=AD-CB$ is one of the following:
\[
 0,\qquad \kappa x_i(x_i-1),\qquad
 \kappa(x_i-x_j)(x_i+x_j-1),
 \tag{1}
\]
where $i\neq j$ and $\kappa\neq0$ in the nonzero cases.
\end{theorem}
\begin{proof}
A polynomial of degree at most two vanishing at every Boolean vertex
has the form
\[
 q=\sum_i\kappa_i(x_i^2-x_i).
 \tag{2}
\]
Indeed reduction modulo $x_i^2=x_i$ gives a multilinear polynomial
vanishing on the cube, whose coefficients are all zero by successive
finite differences at the zero vertex.

Homogenize using coordinate $X_0$. With the affine coefficient vectors
$a,b,c,d$, the symmetric matrix of $q$ is
\[
 Q=\tfrac12(ad^\top+da^\top-cb^\top-bc^\top).
\]
It has rank at most four. Moreover, if $\operatorname{rank}Q\geq3$,
then all four coefficient vectors belong to $\operatorname{im}Q$.
To check this last assertion, let $U$ have columns $a,d,c,b$. Then
$Q=UKU^\top$, where $K$ is an invertible $4$ by $4$ matrix with two
nonzero off-diagonal $2$ by $2$ blocks. If $U$ has rank four, the
surjectivity of $U^\top$ and injectivity of $U$ imply that $Q$ has
rank four and image $\operatorname{im}U$. If $Q$ has rank three,
$U$ cannot have rank four, and consequently both ranks equal three
and their images again agree.

Let $I=\{i:\kappa_i\neq0\}$ and $k=|I|$. Formula (2) gives the
nonzero block of $Q$, indexed by $\{0\}\cup I$, as
\[
 \begin{pmatrix}0&-\kappa^\top/2\\
                 -\kappa/2&\operatorname{diag}(\kappa)\end{pmatrix}.
\]
Its Schur complement with respect to the invertible diagonal block is
$-\sum_i\kappa_i/4$. Hence its rank is $k+1$ if that sum is nonzero,
and $k$ otherwise. In particular $k\leq4$.

Suppose $\operatorname{rank}Q\geq3$. If $k<n$, the image of $Q$
has zero coordinate in every variable outside $I$. The preceding
image assertion forces both denominators to have zero slopes in those
variables. Thus neither denominator can be of the nonconstant type
allowed in the hypothesis. Both would be constant, making $AD-CB$
affine and hence zero, a contradiction.

It remains that $k=n$. Since $n\geq4$ and $k\leq4$, this means
$k=n=4$. The rank bound forces $\sum_i\kappa_i=0$. The homogeneous
cube-center vector $v=(1,1/2,\ldots,1/2)$ then satisfies $Qv=0$.
Symmetry of $Q$ and $b,d\in\operatorname{im}Q$ force
$B(1/2,\ldots,1/2)=D(1/2,\ldots,1/2)=0$. Each value is the average
of the corresponding positive vertex values, a contradiction.

We have proved $\operatorname{rank}Q\leq2$. If $q\neq0$, the rank
formula leaves either $k=1$, or $k=2$ with the two coefficients
opposite. These are exactly the two nonzero factorizations in (1).
\end{proof}

\paragraph{Canonical meaning of the nonzero cases.}
Write $q=\kappa LM$ using its two displayed affine factors.
The denominators $B,D$ are linearly independent. If they were
proportional and constant, $q$ would be affine, which is impossible
in the nonzero cases. If they were proportional and nonconstant,
their common positive affine factor would divide $LM$ and hence
be proportional to $L$ or $M$, but each factor vanishes at a Boolean
vertex.
On the homogeneous subspace $B=D=0$, the product $LM$ vanishes.
A vector space over the infinite field $\mathbb R$ cannot be a union
of two proper linear subspaces. Thus one of the factors, say $L$,
belongs to the span of $B,D$. Write
\[
 D=\alpha B+\beta L.
\]
Here $\beta\neq0$, and evaluating at a Boolean zero of $L$ shows
$\alpha>0$. Since $B$ and $L$ are coprime linear forms, the identity
\[
 (\alpha A-C)B=L(\kappa M-\beta A)
\]
implies $\kappa M-\beta A=k_0B$ for a scalar $k_0$. Therefore
\[
 A=-\frac{k_0}{\beta}B+\frac\kappa\beta M,
 \qquad
 C=-\frac{k_0}{\beta}D+\frac{\alpha\kappa}{\beta}M.
\]
After subtracting a common constant score and positively rescaling
$C,D$, the two numerators coincide and the denominators differ by a
multiple of the other Boolean-vanishing factor. This identifies the
precise one-coordinate or two-coordinate cancellation mechanism.

\paragraph{Exact examples and the dimension threshold.}
For any $n\geq4$, take $B=10+2\sum_i x_i$.
With $D=B+x_1$ and $A=C=x_1-1$, the cross-product is
$x_1(x_1-1)$. With $D=B+x_1-x_2$ and $A=C=x_1+x_2-1$, it is
$(x_1-x_2)(x_1+x_2-1)$. All slopes of both denominators are strictly positive.
Both identities hold on the whole cube and give nonconstant scores.

Dimension three has additional identities, so it cannot be silently
included in the theorem. For variables $x,y,z$, put
\[
 \begin{gathered}
 B=(2885x+5y+6z+184)/5100,\\
 L=1195x+5y+76,\qquad M=1393x-y+92,\\
 K=102(2885x+5y-6z+190),\\
 D=B+L,\qquad A=-M/10,\qquad C=A-K.
 \end{gathered}
\]
All denominator constants and slopes are strictly positive. Direct
expansion gives
\[
 AD-CB=BK-LM/10
      =x(x-1)+y(y-1)-\frac{18}{25}z(z-1).
\]
Thus the two scores agree at all eight vertices while the cross-product
uses three coordinates. The checker verifies this identity coefficient
by coefficient, checks every vertex, and verifies the one-coordinate
and two-coordinate examples in dimension five.
\end{document}
